\documentclass[11pt]{article}
\usepackage[a4paper,margin=1in]{geometry}
\usepackage{fontspec}
\usepackage[boldfont=false]{xeCJK}
\setCJKmainfont{FandolSong-Regular.otf}
\usepackage{amsmath}
\usepackage{amssymb}
\usepackage{array}
\usepackage{booktabs}
\usepackage{graphicx}
\usepackage{adjustbox}
\usepackage{enumitem}
\setlist[itemize]{topsep=2pt,partopsep=0pt,parsep=0pt,itemsep=2pt,leftmargin=1.4em}
\setlist[enumerate]{topsep=2pt,partopsep=0pt,parsep=0pt,itemsep=2pt,leftmargin=1.6em}
\usepackage{parskip}
\usepackage{fvextra}
\fvset{breaklines=true,breakanywhere=true,fontsize=\small}
\setlength{\parindent}{0pt}

\begin{document}

\begin{center}
  {\LARGE\bfseries Halogen compounds}\\[0.35em]
  {\large A Level Chemistry}
\end{center}
\vspace{0.6em}

\section*{Making halogenoarenes}

A \textbf{halogenoarene} 卤代芳烃 (also called an aryl halide) is formed when an \textbf{arene} 芳烃 reacts with $\text{Cl}_2$ or $\text{Br}_2$, using $\text{AlCl}_3$ or $\text{AlBr}_3$ as a \textbf{catalyst} 催化剂. This is the electrophilic substitution from the arenes topic — the halogen replaces a hydrogen on the ring.

\begin{itemize}
\item benzene gives chlorobenzene.

\item methylbenzene gives 2-chloromethylbenzene and 4-chloromethylbenzene (the methyl group directs the chlorine to positions 2 and 4).

\end{itemize}

\section*{Why a halogenoarene is less reactive than a halogenoalkane}

Compare chloroethane (a \textbf{halogenoalkane} 卤代烷) with chlorobenzene (a halogenoarene).

A halogenoalkane reacts easily by \textbf{nucleophilic substitution} 亲核取代: its C–Cl bond is polar, so a nucleophile can attack the slightly positive carbon and push the halogen out.

A halogenoarene is \textbf{very unreactive} towards nucleophilic substitution. There are two reasons:

\begin{itemize}
\item a \textbf{lone pair} 孤对电子 on the chlorine overlaps sideways with the \textbf{delocalised} 离域 ring of electrons. This gives the C–Cl bond partial double-bond character, making it shorter and stronger, so it is much harder to break.

\item the electron-rich ring repels an approaching \textbf{nucleophile} 亲核试剂.

\end{itemize}

So chlorobenzene does \textbf{not} react with nucleophiles such as $\text{OH}^-$ under normal conditions, while chloroethane does.


\end{document}
